
\setcounter{page}{184}
\section{The functor \(f^!\) for embeddable morphisms}

In this section we use the fundamental local isomorphism to relate the functors \(f^*\) and \(f^b\) defined earlier, and define \(f^!\) for morphisms that factor as a finite morphism followed by a smooth morphism. The main results here are provisional, serving as a model for the stronger general duality theorems in Chapter VII after developing local residue techniques.

\begin{lemma}
Let \(f\colon X\to Y\), \(i\colon Y\to X\) be morphisms of locally noetherian preschemes, with \(i\) a section of \(f\) and smooth. Then there exists a functorial isomorphism
\[
\psi_{i,f}\colon G^\bullet \stackrel{\sim}{\to} i^b f^* G^\bullet
\]
for all \(G^\bullet\in D(Y)\).
\end{lemma}

\begin{proof}
By Proposition 1.2, a smooth morphism is locally a complete intersection, so we may apply Corollary 7.3 to \(i\). By definition of \(f^*\),
\[
i^b f^* G^\bullet = i^b\big(\Lfst G^\bullet \otimes \omega_{X/Y}[\dim]\big).
\]
By the fundamental local isomorphism (Corollary 7.3), this is isomorphic to
\[
\Lfst i^*\big(\Lfst G^\bullet \otimes \omega_{X/Y}\big) \otimes \omega_{Y/X}[-\dim].
\]
Using standard derived functor compatibilities [II.5.9], [II.5.4], this simplifies to
\[
G^\bullet \otimes i^*\omega_{X/Y} \otimes \omega_{Y/X}.
\]
Composing with the canonical isomorphism \(\zeta_{i,f}\) from Definition 1.5 yields the isomorphism \(G^\bullet\cong i^b f^* G^\bullet\).
\end{proof}

\setcounter{page}{185}
\begin{proposition}
Let \(f\colon X\to Y\), \(g\colon Y\to Z\) be morphisms of locally noetherian preschemes, with \(f\) finite, \(g\) smooth, and \(gf\) finite. Then there is a functorial isomorphism
\[
\psi_{f,g}\colon (gf)^b \stackrel{\sim}{\to} f^b g^*
\]
defined on \(D_{\mathrm{qc}}^+(Z)\).
\end{proposition}

\begin{proof}
Consider the fibre product \(X\times_Z Y\) with projections \(p_1\colon X\times_Z Y\to X\), \(p_2\colon X\times_Z Y\to Y\), and let \(i\colon X\to X\times_Z Y\) be the graph morphism of \(f\). Since \(g\) is smooth, \(p_1\) is smooth by base change, so Lemma 8.1 applies to \(i,p_1\). We have
\[
(gf)^b \cong i^b p_1^*(gf)^b.
\]
By Corollary 6.4 this is isomorphic to \(i^b p_2^* g^*\), and by Proposition 6.2 this reduces to \(f^b g^*\).
\end{proof}

\setcounter{page}{186}
\begin{remark}
This residue isomorphism was first discovered by Grothendieck via a lengthy argument. The streamlined proof above is due to Cartier (as interpreted by Mumford). It plays a critical role in constructing the trace map for residual complexes in Chapter VI, a key prerequisite for the general residue theorem.
\end{remark}

\begin{corollary}
Let \(f\colon X\to Y\) be a morphism of locally noetherian preschemes which is both finite and smooth. Then there is an isomorphism
\[
\psi_f\colon f^b \stackrel{\sim}{\to} f^*
\]
defined on \(D_{\mathrm{qc}}^+(Y)\).
\end{corollary}

\begin{proof}
Set \(g=\id_Y\) in Proposition 8.2.
\end{proof}

\setcounter{page}{187}
\begin{remark}
One may verify that this isomorphism coincides with the classical trace map \(f_*\sheaf{O}_X\to\sheaf{Y}\) when restricted to ordinary sheaves.
\end{remark}

\begin{proposition}
Let \(f\colon X\to Y\), \(g\colon Y\to Z\) be morphisms of locally noetherian preschemes, with \(f\) finite, \(g\) smooth, and \(gf\) smooth. Then there is a functorial isomorphism
\[
\psi_{f,g}\colon (gf)^* \to f^b g^*
\]
defined on \(D_{\mathrm{qc}}^+(Z)\).
\end{proposition}

\begin{proof}
Using the fibre product notation from Proposition 8.2:
\[
(gf)^* \cong i^b p_1^*(gf)^*
\cong i^b p_2^* g^*
\cong f^b g^*.
\]
The isomorphisms follow respectively from Lemma 8.1, Proposition 2.2, and Proposition 8.2.
\end{proposition}

\begin{corollary}
With the hypotheses of Proposition 8.4, there exists a natural trace morphism
\[
\operatorname{Tr}_f\colon f_*\omega_{X/Z} \to \omega_{Y/Z}.
\]
\end{corollary}

\begin{proof}
Apply the isomorphism of Proposition 8.4 to \(G^\bullet=\sheaf{O}_Z\) and compose with the trace map from Proposition 6.5.
\end{proof}

\setcounter{page}{188}
\begin{remark}
When \(Z=\operatorname{Spec}k\) for a field \(k\), and \(X,Y\) are irreducible with separable function field extension \(K(X)/K(Y)\), this trace morphism recovers the classical field trace. It satisfies standard functorial compatibilities with composition and flat base extension, and merits further independent study.
\end{remark}

\begin{proposition}
\begin{enumerate}
\item[(a)] The isomorphisms of Propositions 8.2 and 8.4 are compatible with flat base extension, via the base-change isomorphisms of Propositions 2.1 and 6.3.
\item[(b)] For three composable morphisms \(X\xrightarrow{f}Y\xrightarrow{g}Z\xrightarrow{h}W\) satisfying the hypotheses of Propositions 2.2, 6.2, 8.2, 8.4 pairwise, the corresponding functor isomorphisms form commutative diagrams.
\item[(c)] For Cartesian diagrams with finite/smooth morphisms, the induced isomorphisms are canonically compatible.
\end{enumerate}
\end{proposition}